\section{Appendix: Scope and Qualifications}
This study concerns All Saints in the Vatican typical edition of the
1962 Roman Missal, under the General Roman Calendar, on 1 November 2026.
The privileged commemoration of the Twenty-third Sunday after Pentecost
adds its three orations. No diocesan, religious, titular, or patronal
calendar is supplied. The distinct votive Mass of All Saints and its
seasonal variants are excluded. No present canonical permission to use
the book, local schedule, or territorial holy-day discipline is determined.

The complete Latin proper bodies, including the Common Preface body,
come from the research's target collation and identified public-domain
antecedents. The stable Preface dialogue and Sanctus belong to the
Ordinary, outside the proper body printed here. Abbreviated conclusions
remain abbreviated. The independent Douay--Rheims/Challoner study
selections retain their biblical boundaries; historical Lasance/Walsh
English identifies the Missal's added framing and adaptations.
Cummiskey supplies historical English for the Common Preface. None is
presented as approved English for the 1962 Missal. The original
commentary is distinct from every received prayer.

The two interpretations are editorial syntheses of this Mass. Bede,
Augustine, Chrysostom, Aquinas, and Gregory comment on the passages or
doctrinal questions specified in their citations. Schuster supplies
reception of the ten feast texts and the three shared Sunday prayers;
his older rubrics and historical Preface do not determine this dated
celebration. English transcriptions supply the stated patristic and
medieval arguments; no fresh Greek or critical Latin collation is claimed.
The bounded source sweep, negative results, text collation, chronology,
and reusable-source identities are recorded in \texttt{research/scope.md},
\texttt{research/interpretations.md}, \texttt{research/source-bindings.toml},
and \texttt{propers/verified.md}.

The Latin and historical English retain their recorded United States
public-domain bases; the project claims no exclusive right in them.
The modern 1962 facsimile, modern scholarly introductions, and the
Corpus Thomisticum edition are inspection sources, not redistributed
whole-source payloads. The research's Gloria observation remains bounded
to the hand-missal direction and inspected Missal rule; this study adds
no claim of independent verification of its Breviary reference.
The author's proof checks are recorded separately from the independent
reviews required for the three completed documents.

\section{References}
\subsection{Liturgical and biblical sources}
\begin{itemize}
\item \work{Missale Romanum}, Vatican editio typica (1962),
All Saints, pp.~702--703, nos.~3974--3983; commemoration prayers,
pp.~416--417, nos.~1720, 1726, 1728; Common Preface,
pp.~252--253, no.~1057. \work{Rubricae generales} 91, 106,
108--109, 111, 113, 120(d); \work{Rubricae generales Missalis romani}
475, 480, 482--483, 494(b), 498, 505. Controlling
\href{https://media.churchmusicassociation.org/pdf/missale62.pdf}{CMAA facsimile}.
\item Public-domain Latin antecedents: \work{Missale Romanum},
Mame, Tours (1922), pp.~748--750, All Saints;
Pustet, Ratisbon (1862), pp.~357--358, Sunday orations;
Eugene Cummiskey, \work{The Roman Missal translated into the
English language for the use of the laity}, first revised edition
(Philadelphia, 1861), p.~xxviii, Common Preface Latin and English.
\item F.~X. Lasance and F.~A. Walsh, \work{The New Roman Missal},
revised 1945, ordinary-numbered pp.~1265--1268, All Saints;
pp.~733, 735--736, three Sunday orations. Historical human
English; later lettered inserts are excluded.
\item \work{Douay--Rheims Bible}, Challoner revision,
\href{https://www.gutenberg.org/ebooks/1581}{Gutenberg 1581}:
Ps.~32:1; 33:10--11; Wis.~2:12--3:9; Matt.~4:23--5:12;
11:25--30; Apoc.~1:4,9; 6:12--7:17; Gal.~3:29.
Vulgate Psalm numbering is retained.
\end{itemize}

\subsection{Commentaries and theological sources}
\begin{itemize}
\item Bede, \work{The Explanation of the Apocalypse}, trans.
Edward Marshall (Oxford and London: James Parker, 1878),
book I on Apoc.~7:4--12, pp.~46--53; context 7:13--17,
pp.~53--54. The registered complete facsimile and bounded checked
comments identify the historical English witness.
\item Augustine, \work{Our Lord's Sermon on the Mount},
I.1.3--I.4.12, NPNF first series 6, William Findlay's historical
English with Schaff's revision and editorial notes; retained
CCEL text. Notes by the editor remain distinct from Augustine.
\item Augustine, \work{Expositions of the Psalms}, Ps.~32,
enarratio I, \S1; Ps.~33, enarratio II, \S\S14--15;
Ps.~87, \S5. NPNF first series 8, retained CCEL historical
English: Psalm XXXIII \S1, Psalm XXXIV \S\S13--14,
and Psalm LXXXVIII \S4 respectively.
\item Augustine, \work{The City of God}, XX.9, Dods edition,
II (1871), retained Gutenberg historical English. Its direct
Apocalypse discussion concerns chapter 20.
\item John Chrysostom, \work{Homilies on Matthew}, XV.2--7
and XXXVIII.3--5, NPNF first series 10, historical English
revised by George Prevost and M.~B. Riddle; retained CCEL text.
\item Thomas Aquinas, \work{Super Psalmos}, Psalm 33,
nn.~10--11, \href{https://www.corpusthomisticum.org/cps31.html}{
Corpus Thomisticum Latin edition}; paraphrased, without
reproducing its modern editorial text.
\item Gregory the Great, \work{Moralia in Iob}, XII.5.6--7,
on Job~14:8--9, directly reusing Wis.~3:2; historical
\work{Library of the Fathers} English in the registered
complete Lectionary Central transcription.
\item Alfredo Ildefonso Schuster, \work{The Sacramentary},
III (London: Burns Oates \& Washbourne, 1927), pp.~186--189,
three Sunday orations; V (1930), pp.~208--211, All Saints.
The historical pages control the cited paraphrases.
\end{itemize}

The historical appendix cites the particular \work{Catholic
Encyclopedia} articles, USCCB introduction, and Haydock
facsimile loci used for its chronological and geographical claims.
